\documentclass[11pt]{article}
\usepackage[margin=1in]{geometry}
\usepackage{amsmath,amssymb}
\usepackage[hidelinks]{hyperref}
\usepackage{parskip}

\begin{document}

\begin{flushright}
Sparsho Chakraborty \\
School of Basic Sciences, Department of Physics, \\
Indian Institute of Technology Bhubaneswar, Odisha, India \\
sparshochakraborty123@gmail.com \\
ORCID: 0009-0004-1667-7208
\end{flushright}

\bigskip
\noindent \textbf{To:} Editor-in-Chief, \textit{Quantum Machine Intelligence} (Springer)

\bigskip
\noindent \textbf{Re:} Submission of the manuscript ``Hamiltonian Vision Kernel: Symmetry-Pooled Quantum Correlator Features with Hardware Validation''

\bigskip
\noindent Dear Editor,

\smallskip
\noindent We are submitting the enclosed manuscript, ``Hamiltonian Vision Kernel: Symmetry-Pooled Quantum Correlator Features with Hardware Validation,'' for consideration as a research article in \textit{Quantum Machine Intelligence}.

\smallskip
\noindent The manuscript introduces the Hamiltonian Vision Kernel (HVK), a new hybrid quantum-classical architecture for image reconstruction that combines matrix-product-state patch features, a variational quantum circuit exposing a Pauli-correlator latent space, Fourier positional encoding, a learnable Hamiltonian energy diagnostic, and a classical decoder. We make no claim of quantum advantage: HVK is presented as a new, working architecture that is competitive with resource-matched classical baselines, validated on real quantum hardware, and equipped with capabilities --- an entanglement-sensitive channel and an interpretable Hamiltonian diagnostic --- that a classical reconstruction pipeline does not natively expose. We believe this fits the scope of \textit{Quantum Machine Intelligence} directly: it sits at the intersection of quantum machine learning, hybrid classical-quantum architectures, and tensor-network methods for structured data, evaluated on real quantum hardware, and it takes deliberate care to characterize where the approach's evidence is strong and where it is exploratory rather than confirmatory.

\smallskip
\noindent Five results anchor the manuscript:

\begin{enumerate}
\item The pooled HVK2D observable map achieves $D_4$ group-equivariance by construction, with a measured numerical equivariance error of $9.57 \times 10^{-17}$.
\item On a resource-matched, multi-seed held-out CIFAR-10 comparison, a pre-declared Two One-Sided Tests (TOST) equivalence test at a $\pm1$~dB margin shows HVK2D is statistically equivalent to six of eight tested classical/ablated controls, and substantially outperforms the two null controls (random-VQC, strict classical random features) --- the evidence that licenses calling HVK ``competitive'' rather than simply ``not shown to differ.''
\item On a leakage-audited diagnostic task requiring distant-patch correlations, the entangling observable channel reaches mean $R^2 = 0.974$ across five seeds, against $R^2 \leq 0.02$ for non-entangling controls -- evidence that the model's entangling structure, not incidental capacity, is responsible for this capability.
\item Five images are reconstructed from observables measured on real IBM Quantum hardware, without retraining the decoder, reaching 25.90--31.52 dB PSNR -- a scoped hardware feasibility pilot, reported with complete job metadata and simulator-transfer checks rather than as a general performance claim.
\item We report the boundary of the approach as plainly as its strengths. On ordinary held-out reconstruction the simple classical controls are numerically ahead of HVK ($18.80$ vs.\ $18.12$~dB), and we say so rather than burying it; the learned energy and entanglement are reported only as an interpretability readout, together with the negative control that disqualifies any stronger reading of them --- a threshold applied to those traces fires just as often on a classical-replacement variant whose Hamiltonian energy is identically zero ($4/4$ vs.\ $4/4$). We claim no transition, no critical epoch, and no critical temperature anywhere in the manuscript.
\end{enumerate}

\smallskip
\noindent A companion supplementary study accompanies the manuscript with resource-matched held-out controls, leakage audits, statistical equivalence testing, and complete hardware execution methodology (backend names, job IDs, shot counts). Code, trained checkpoints, and raw result artifacts backing every reported number are maintained in a public repository (see \texttt{REPRODUCE.md} in the submission materials for a one-to-one table/figure-to-script map).

\smallskip
\noindent This manuscript is not under review at any other journal, and all authors have approved this submission. We intend to post a preprint on arXiv, consistent with Springer Nature's self-archiving policy, and are submitting via the standard subscription route.

\bigskip
\noindent \emph{[FILL: Suggested reviewers (2--3 names, affiliations, emails). Needs domain knowledge of who is active and appropriate in variational quantum algorithms / quantum machine learning / tensor-network methods for images -- a judgment call for the authors and supervisor, not something to guess.]}

\bigskip
\noindent We thank the editors and reviewers for their time and consideration.

\bigskip
\noindent Sincerely,

\bigskip
\noindent Sparsho Chakraborty \\
School of Basic Sciences, Department of Physics, Indian Institute of Technology Bhubaneswar, Odisha, India \\
sparshochakraborty123@gmail.com, 25ph05023@iitbbs.ac.in \\
ORCID: 0009-0004-1667-7208

\bigskip
\noindent Siddhartha Patra \\
Assistant Professor, Centre for Quantum Engineering, Research and Education (CQuERE), TCG CREST, Kolkata, India \\
siddhartha.patra@tcgcrest.org \\
ORCID: \emph{[FILL]}
\end{document}
